% part: intuitionistic-logic
% chapter: introduction
% section: bhk-interpretation

\documentclass[../../../include/open-logic-section]{subfiles}

\begin{document}

\olfileid{int}{int}{bhk}

\olsection{Interpretasi Brouwer--Heyting--Kolmogorov}

\begin{editorial}
  Bukti validitas proposisi intuisionistik nganggo interpretasi BHK
  isih mbingungake; perlu dijlentrehake kanthi luwih cetha.
\end{editorial}

Ana interpretasi konstruktif informal kanggo konektif intuisionistik,
kang lumrahe diarani interpretasi Brouwer--Heyting--Kolmogorov.
Interpretasi iki nggunakake gagasan ``konstruksi,'' kang bisa dianggep
minangka bukti konstruktif. (Ing interpretasi BHK kita ora nggunakake
tembung ``bukti'' supaya ora kacampur karo !!a{derivation} ing sistem
!!{derivation} formal.) Adhedhasar gagasan intuisi iki, interpretasi
BHK njlentrehake teges konektif intuisionistik.

\begin{enumerate}
\item Kita nganggep wis ngerti apa kang diarani konstruksi kanggo
  pratelan atomik.
\item Konstruksi kanggo $!A_1 \land !A_2$ yaiku pasangan $\tuple{M_1, M_2}$
  kang $M_1$ iku konstruksi kanggo~$!A_1$ lan $M_2$ iku konstruksi
  kanggo~$!A_2$.
\item Konstruksi kanggo $!A_1 \lor !A_2$ yaiku pasangan $\tuple{s, M}$
  kang $s$ iku~$1$ lan $M$ iku konstruksi kanggo~$!A_1$, utawa $s$ iku~$2$
  lan $M$ iku konstruksi kanggo~$!A_2$.
\item Konstruksi kanggo $!A \lif !B$ yaiku fungsi kang ngowahi
  konstruksi kanggo~$!A$ dadi konstruksi kanggo~$!B$.
\item Ora ana konstruksi kanggo~$\lfalse$ (absurditas).
\item $\lnot !A$ ditetepake minangka sinonim saka $!A \lif \lfalse$.
  Tegese, konstruksi kanggo $\lnot !A$ yaiku fungsi kang ngowahi
  konstruksi kanggo~$!A$ dadi konstruksi kanggo~$\lfalse$.
\end{enumerate}

\begin{ex}
Contone $\lnot \lfalse$. Konstruksi kanggo rumus iki yaiku fungsi
kang nampa sembarang konstruksi kanggo $\lfalse$ minangka input lan
menehi konstruksi kanggo~$\lfalse$ minangka output. Temenan, fungsi
identitas~$\Id{}$ iku konstruksi kaya mangkono: yen diwenehi
konstruksi~$M$ kanggo~$\lfalse$, $\Id{}(M) = M$ ngasilake
konstruksi kanggo~$\lfalse$.
\end{ex}

Sacara umum, $\lnot !A$ ateges ``Konstruksi kanggo $!A$ iku mokal''.

\begin{ex}
Saiki ayo buktekake $!A \lif \lnot \lnot !A$ kanggo saben
proposisi~$!A$, yaiku $!A \lif ((!A \lif \lfalse) \lif \lfalse)$.
Konstruksine kudu awujud fungsi~$f$ kang, yen diwenehi konstruksi~$M$
kanggo $!A$, ngasilake konstruksi~$f(M)$ kanggo
$(!A \lif \lfalse) \lif \lfalse$. Mangkene carane~$f$ mbangun
konstruksi kanggo $(!A \lif \lfalse) \lif
\lfalse$: kita kudu netepake fungsi $g$ kang nampa konstruksi~$h$
kanggo $!A \lif \lfalse$ minangka input lan ngasilake konstruksi
kanggo~$\lfalse$. Fungsi $g$ bisa ditetepake kanthi ngetrapake
input~$h$ marang konstruksi~$M$ kanggo $!A$ kang wis ditampa.
Amarga output~$h(M)$ saka $h$ iku konstruksi kanggo $\lfalse$,
$f(M)(h) = h(M)$ iku konstruksi kanggo~$\lfalse$ yen $M$ iku
konstruksi kanggo~$!A$.
\end{ex}

\begin{ex}
Saiki kita wenehi konstruksi kanggo $\lnot(!A \land \lnot !A)$,
yaiku $(!A
\land (!A \lif \lfalse)) \lif \lfalse$. Iki fungsi~$f$ kang,
yen nampa konstruksi~$M$ kanggo $!A \land (!A \lif \lfalse)$,
ngasilake konstruksi kanggo~$\lfalse$. Konstruksi kanggo konjungsi
$!B_1 \land !B_2$ yaiku pasangan $\tuple{N_1, N_2}$ kang $N_1$
iku konstruksi kanggo~$!B_1$ lan $N_2$ iku konstruksi kanggo~$!B_2$.
Kita bisa netepake fungsi $p_1$ lan $p_2$ kang njupuk saka konstruksi
kanggo $!B_1 \land !B_2$ konstruksi kanggo $!B_1$ lan $!B_2$
miturut urutane:
\begin{align*}
  p_1(\tuple{N_1, N_2}) &= N_1\\
  p_2(\tuple{N_1, N_2}) &= N_2
\end{align*}
Mangkene tumindake $f$: dhisik fungsi iku ngetrapake $p_1$ marang
input~$M$, mula entuk konstruksi kanggo~$!A$. Banjur $p_2$
ditrapake marang~$M$ lan ngasilake konstruksi kanggo
$!A \lif \lfalse$. Konstruksi iki yaiku fungsi~$p_2(M)$ kang,
yen nampa konstruksi kanggo~$!A$, ngasilake konstruksi
kanggo~$\lfalse$. Dadi, yen $p_2(M)$ ditrapake marang $p_1(M)$,
kita entuk konstruksi kanggo~$\lfalse$. Mula bisa ditetepake
$f(M) = p_2(M)(p_1(M))$.
\end{ex}

\begin{ex}
Ayo menehi konstruksi kanggo $((!A \land !B) \lif !C) \lif (!A \lif
(!B \lif !C))$, yaiku fungsi $f$ kang ngowahi konstruksi~$g$ kanggo
$(!A \land !B) \lif !C$ dadi konstruksi kanggo $(!A \lif (!B \lif
!C))$. Konstruksi $g$ iku dhewe fungsi saka konstruksi kanggo
$!A \land !B$ menyang konstruksi kanggo~$!C$. Output $f(g)$ yaiku
fungsi~$h_g$ saka konstruksi kanggo $!A$ menyang fungsi kang ngowahi
konstruksi kanggo~$!B$ dadi konstruksi kanggo~$!C$.

Ya, iki pancen mbingungake. Kita kudu mbangun fungsi $h_g$ kang dadi
output saka $f$ kanggo input~$g$. Input kanggo~$h_g$ yaiku
konstruksi~$M$ kanggo~$!A$. Output $h_g(M)$ kudu awujud fungsi~$k_{g,M}$
kang ngowahi konstruksi~$N$ kanggo~$!B$ dadi konstruksi kanggo~$!C$.
Netepna $k_{g,M}(N) = g(\tuple{M, N})$. Elinga yen $\tuple{M,
  N}$ iku konstruksi kanggo~$!A \land !B$. Dadi $k_{g,M}$ iku
konstruksi kanggo $!B \lif !C$: fungsi iki ngowahi konstruksi~$N$
kanggo~$!B$ dadi konstruksi kanggo~$!C$. Saiki netepna
$h_g(M) = k_{g,M}$. Iki fungsi kang ngowahi konstruksi~$M$ kanggo~$!A$
dadi konstruksi $k_{g,M}$ kanggo $!B
\lif !C$. Banjur netepna $f(g) = h_g$. Iki fungsi kang ngowahi
konstruksi $g$ kanggo $(!A \land !B) \lif !C$ dadi konstruksi kanggo
$!A
\lif (!B \lif !C)$. Rampung!{}
\end{ex}

Pratelan $!A \lor \lnot !A$ diarani Hukum Tengah kang Disingkirake.
Kita bisa mbuktekake kanggo $!A$ tartamtu (umpamane $\lfalse \lor
\lnot \lfalse$), nanging ora kanthi umum. Sebabe, disjungsi
intuisionistik mbutuhake konstruksi kanggo salah siji cabange, dene
ana pratelan kang saiki durung bisa dibuktekake utawa dibantah
(umpamane dugaan Goldbach). Nanging Hukum Tengah kang Disingkirake uga
ora bisa dibantah: yaiku, $\lnot \lnot (!A
\lor \lnot !A)$ bener.

\begin{ex}
Kanggo mbuktekake $\lnot \lnot (!A \lor \lnot !A)$, kita perlu
fungsi~$f$ kang ngowahi konstruksi kanggo $\lnot(!A \lor \lnot !A)$,
yaiku kanggo $(!A
\lor (!A \lif \lfalse)) \lif \lfalse$, dadi konstruksi
kanggo~$\lfalse$. Kanthi tembung liya, kita perlu fungsi $f$ kang
$f(g)$ iku konstruksi kanggo~$\lfalse$ yen $g$ iku konstruksi kanggo
$\lnot(!A
\lor \lnot !A)$.

Upamane $g$ iku konstruksi kanggo $\lnot(!A \lor \lnot !A)$, yaiku
fungsi kang ngowahi konstruksi kanggo $!A \lor \lnot !A$ dadi
konstruksi kanggo~$\lfalse$. Konstruksi kanggo $!A \lor \lnot !A$
iku pasangan $\tuple{s, M}$ kang $s = 1$ lan $M$ iku konstruksi kanggo
$!A$, utawa $s = 2$ lan $M$ iku konstruksi kanggo $\lnot !A$.
Upamane $h_1$ iku fungsi kang nggawa konstruksi~$M_1$ kanggo $!A$
menyang konstruksi kanggo $!A \lor \lnot !A$: fungsi iku nggawa
$M_1$ menyang $\tuple{1, M_1}$. Dene $h_2$ iku fungsi kang nggawa
konstruksi~$M_2$ kanggo $\lnot !A$ menyang konstruksi kanggo
$!A \lor \lnot !A$: fungsi iku nggawa $M_2$ menyang $\tuple{2, M_2}$.

Upamane $k$ iku $\comp{h_1}{g}$: fungsi iki, yen diwenehi konstruksi
kanggo~$!A$, ngasilake konstruksi kanggo~$\lfalse$. Dadi, iki
konstruksi kanggo~$!A \lif \lfalse$ utawa $\lnot !A$. Saiki upamane
$l$ iku $\comp{h_2}{g}$. Fungsi iki, yen diwenehi konstruksi kanggo
$\lnot !A$, ngasilake konstruksi kanggo~$\lfalse$. Amarga $k$ iku
konstruksi kanggo $\lnot !A$, $l(k)$ iku konstruksi kanggo~$\lfalse$.

Kanthi mangkono, kita wis nerangake carane ngowahi konstruksi~$g$
kanggo $\lnot(!A \lor \lnot !A)$ dadi konstruksi kanggo~$\lfalse$.
Tegese, fungsi $f$ kang nggawa konstruksi~$g$ kanggo
$\lnot(!A \lor \lnot !A)$ menyang konstruksi $l(k)$ kanggo~$\lfalse$
iku konstruksi kanggo $\lnot\lnot(!A \lor \lnot !A)$.
\end{ex}

Kaya kang katon, nggunakake interpretasi BHK kanggo nuduhake
validitas intuisionistik !!{formula}s enggal dadi ruwet lan
mbingungake. Untunge, ana sistem !!{derivation} kang luwih becik
kanggo logika intuisionistik, uga interpretasi semantik kang luwih tliti.
\end{document}
